\section{Longest Increasing Subsequence using Binary Search}
\label{sec:dp-lis-binary-search}

\problemstatement{Solve Section~\ref{sec:dp-lis} (find the length of the
LIS) in $O(n \log n)$ instead of $O(n^2)$.}

\begin{patternbox}
\emph{``Find the length of the LIS, faster than $O(n^2)$''} is solved
not by a better DP table, but by an entirely different technique:
maintain a single array $\textit{tails}$, where $\textit{tails}[k]$ is
the \textbf{smallest possible value} that can end an increasing
subsequence of length $k+1$, discovered so far -- and update it for
each new element using \textbf{binary search} (the Binary Search chapter's
lower bound) instead of scanning all previous elements.
\end{patternbox}

\subsection*{Understanding why \textit{tails} is always sorted}

The key invariant: $\textit{tails}$ is always kept sorted in increasing
order. This is not obvious at first, but follows by induction: a longer
increasing subsequence's smallest possible tail value can never be
smaller than a shorter one's smallest possible tail value (if it were,
that shorter prefix of the longer subsequence would itself be a
counterexample to the shorter subsequence's tail already being
minimal). Because $\textit{tails}$ stays sorted, we can binary search it.

\begin{ideabox}[Update Rule for Each New Element]
For each $x = \textit{arr}[i]$, in order: binary search $\textit{tails}$
for the leftmost position $\textit{pos}$ where $\textit{tails}[\textit{pos}]
\ge x$ (this is exactly the Binary Search chapter's lower bound technique).
If no such position exists ($x$ is larger than every current tail),
append $x$ (a new, longer subsequence is now possible). Otherwise,
replace $\textit{tails}[\textit{pos}] \gets x$ ($x$ is a \emph{smaller}
valid tail for a subsequence of that same length, strictly improving
future extension possibilities, without changing the \emph{count} of
distinct lengths discovered).
\end{ideabox}

\subsection*{Why replacing (not just appending) preserves correctness}

$\textit{tails}$'s length, at any point, equals the LIS length found so
far among the elements processed -- but the actual \emph{values} stored
are not necessarily a real subsequence from the array; they are simply
the best (smallest) tail achievable for each length. Replacing
$\textit{tails}[\textit{pos}]$ with a smaller value $x$ never decreases
$\textit{tails}$'s length (we are not removing entries, only improving
one in place), and it can only ever \emph{help} future elements extend a
chain of that length, since a smaller tail is easier for a later,
possibly-not-much-larger element to exceed. The final length of
$\textit{tails}$ is therefore exactly the LIS length.

\subsection*{Algorithm}

\begin{enumerate}
  \item Initialize an empty array \textit{tails}.
  \item For each $x$ in $\textit{arr}$: binary search for the lower
        bound of $x$ in \textit{tails}. If found at the end (no entry
        $\ge x$): append $x$. Else: overwrite that entry with $x$.
  \item Return $|\textit{tails}|$.
\end{enumerate}

\subsection*{Dry run}

\dryrun{Take $\textit{arr} = [10,9,2,5,3,7,101,18]$.}

\begin{itemize}
  \item $x=10$: \textit{tails} empty, append. $\textit{tails}=[10]$.
  \item $x=9$: lower bound of $9$ in $[10]$ is index $0$. Overwrite.
        $\textit{tails}=[9]$.
  \item $x=2$: lower bound of $2$ in $[9]$ is index $0$. Overwrite.
        $\textit{tails}=[2]$.
  \item $x=5$: lower bound of $5$ in $[2]$: none $\ge5$, append.
        $\textit{tails}=[2,5]$.
  \item $x=3$: lower bound of $3$ in $[2,5]$ is index $1$ (value $5$).
        Overwrite. $\textit{tails}=[2,3]$.
  \item $x=7$: lower bound of $7$ in $[2,3]$: none $\ge7$, append.
        $\textit{tails}=[2,3,7]$.
  \item $x=101$: append. $\textit{tails}=[2,3,7,101]$.
  \item $x=18$: lower bound of $18$ in $[2,3,7,101]$ is index $3$
        (value $101$). Overwrite. $\textit{tails}=[2,3,7,18]$.
\end{itemize}

Final length: $4$ -- matching Section~\ref{sec:dp-lis}'s answer
(note $\textit{tails}=[2,3,7,18]$ is not necessarily a real
subsequence of $\textit{arr}$, but its \emph{length} is correctly the
LIS length).

\subsection*{C++ implementation}

\begin{lstlisting}
#include <bits/stdc++.h>
using namespace std;

int lisBinarySearch(vector<int>& arr) {
    vector<int> tails;

    for (int x : arr) {
        int pos = lower_bound(tails.begin(), tails.end(), x) - tails.begin();
        if (pos == (int)tails.size()) {
            tails.push_back(x);
        } else {
            tails[pos] = x;
        }
    }
    return tails.size();
}

int main() {
    vector<int> arr = {10, 9, 2, 5, 3, 7, 101, 18};
    cout << lisBinarySearch(arr) << endl; // 4
    return 0;
}
\end{lstlisting}

\begin{complexitybox}
\textbf{Time Complexity.} Each of the $n$ elements performs one binary
search over \textit{tails} (size at most $n$), giving
\[
  O(n \log n).
\]
\textbf{Space Complexity.} $O(n)$ for the \textit{tails} array.
\end{complexitybox}

\begin{takeawaybox}
Not every problem's fastest solution is a bigger or cleverer DP table --
sometimes, as here, an entirely different paradigm (greedy maintenance
of a sorted auxiliary array, queried by binary search) beats DP's
polynomial bound outright. Keep this $O(n \log n)$ technique available
specifically for when a problem only needs the LIS \emph{length}; it
does not directly reconstruct the actual subsequence the way
Section~\ref{sec:dp-print-lis}'s parent-pointer approach does.
\end{takeawaybox}
